\begin{helper}
Let \(\Omega_0\) be a finite family of paired supported source and target
cells \(C^0_{\omega_0}\) and \(C^{\prime 0}_{\omega_0}\).  Assume both
families are pairwise disjoint.  Let \(\Lambda\) be a finite list of secondary
Boolean tests, with source tests \(Q_\lambda\) and target tests
\(Q'_\lambda\).  Suppose that the source and target base cells are transported
by a relation \(T_{\omega_0}\), that every secondary test is transported by
the same relation, and that for every base index \(\omega_0\) and truth vector
\(\tau\subseteq\Lambda\) the source and target Boolean subcells
\[
C^0_{\omega_0}\cap
  \{\,\eta: \eta\in Q_\lambda \Longleftrightarrow \lambda\in\tau
       \text{ for all }\lambda\in\Lambda\,\}
\]
and
\[
C^{\prime 0}_{\omega_0}\cap
  \{\,\eta': \eta'\in Q'_\lambda \Longleftrightarrow \lambda\in\tau
       \text{ for all }\lambda\in\Lambda\,\}
\]
have equal source and target measure, and suppose this common atom mass is
represented by a nonnegative real weight.  Assume also that, for every
\(\omega_0\), the source base-cell mass and the target base-cell mass already
split over these finitely many Boolean truth-vector atoms:
\[
\nu(C^0_{\omega_0})
=\sum_{\tau\subseteq\Lambda}
  \nu\!\left(C^0_{\omega_0}\cap
  \{\,\eta: \eta\in Q_\lambda \Longleftrightarrow \lambda\in\tau
       \text{ for all }\lambda\,\}\right),
\]
\[
\nu'(C^{\prime0}_{\omega_0})
=\sum_{\tau\subseteq\Lambda}
  \nu'\!\left(C^{\prime0}_{\omega_0}\cap
  \{\,\eta': \eta'\in Q'_\lambda \Longleftrightarrow \lambda\in\tau
       \text{ for all }\lambda\,\}\right).
\]
Then there is one finite refinement
\(\Omega\), with a map \(b:\Omega\to\Omega_0\), a truth vector
\(\tau_\omega\subseteq\Lambda\), source cells \(C_\omega\), target cells
\(C'_\omega\), and nonnegative weights \(w_\omega\), such that
\[
C_\omega=C^0_{b(\omega)}\cap
  \{\,\eta: \eta\in Q_\lambda \Longleftrightarrow \lambda\in\tau_\omega
       \text{ for all }\lambda\,\},
\]
\[
C'_\omega=C^{\prime 0}_{b(\omega)}\cap
  \{\,\eta': \eta'\in Q'_\lambda \Longleftrightarrow \lambda\in\tau_\omega
       \text{ for all }\lambda\,\}.
\]
The refined source cells and refined target cells are pairwise disjoint, they
cover the corresponding base cells fiber-by-fiber over \(\Omega_0\), all
secondary tests are constant on each refined source and target cell with
truth vector \(\tau_\omega\), transported states have equivalent membership in
\(C_\omega\) and \(C'_\omega\), and
\[
\operatorname{ofReal}(w_\omega)=\nu(C_\omega)=\nu'(C'_\omega).
\]
Moreover the base-cell masses split as the finite sums of the refined-cell
masses over each fiber \(b^{-1}(\omega_0)\), on both the source and target
sides.
\end{helper}
\begin{proof}
Enumerate the finite set of truth vectors \(\mathcal P(\Lambda)\).  Take
\(\Omega=\Omega_0\times\mathcal P(\Lambda)\), let \(b(\omega_0,\tau)=\omega_0\),
and let the truth vector attached to \((\omega_0,\tau)\) be \(\tau\).  Define
the source and target refined cells by the two displayed formulas in the
statement.

For a fixed base index \(\omega_0\), every state in \(C^0_{\omega_0}\) has a
unique truth vector, namely the set of labels \(\lambda\) for which the state
lies in \(Q_\lambda\).  Hence the refined source cells over \(\omega_0\) cover
\(C^0_{\omega_0}\).  Two different truth vectors disagree on some label
\(\lambda\), so their atoms are disjoint.  If two refined source cells have
different base indices, they are disjoint by the assumed disjointness of the
base cells.  This proves the source cover and pairwise disjointness; the
target proof is identical with \(C^{\prime0}_{\omega_0}\) and \(Q'_\lambda\).

The definitions make every secondary test constant on each refined cell:
membership in \(Q_\lambda\), respectively \(Q'_\lambda\), is equivalent to
\(\lambda\in\tau\).  If \(T_{\omega_0}(\eta,\eta')\) transports a source state
to a target state, the assumed base-cell transport identifies membership in
the paired base cells, and the assumed label transport identifies every
Boolean label.  Therefore the full conjunction defining \(C_{(\omega_0,\tau)}\)
is true for \(\eta\) exactly when the full conjunction defining
\(C'_{(\omega_0,\tau)}\) is true for \(\eta'\).

By the real atom-mass hypothesis for each Boolean atom, the two refined cells
with the same \((\omega_0,\tau)\) have a common nonnegative real
representative.  Choose \(w_\omega\) to be that representative.  This gives
\(\operatorname{ofReal}(w_\omega)=\nu(C_\omega)=\nu'(C'_\omega)\).

It remains only to verify that the final displayed sums are indexed in the
same way as the assumed base-to-atom mass splittings.  For a fixed
\(\omega_0\), the fiber \(b^{-1}(\omega_0)\) is precisely
\(\{\omega_0\}\times\mathcal P(\Lambda)\).  Therefore
\[
\sum_{\omega\in\Omega}
 {\bf 1}_{b(\omega)=\omega_0}\,\nu(C_\omega)
=\sum_{\tau\subseteq\Lambda}
  \nu(C_{(\omega_0,\tau)}),
\]
and by the definition of \(C_{(\omega_0,\tau)}\) this is exactly the source
truth-vector sum assumed in the statement.  The source mass-splitting
hypothesis identifies that sum with \(\nu(C^0_{\omega_0})\).  The identical
argument on the target side, using the target mass-splitting hypothesis and
the definition of \(C'_{(\omega_0,\tau)}\), gives
\[
\nu'(C^{\prime0}_{\omega_0})
=\sum_{\omega\in\Omega}
 {\bf 1}_{b(\omega)=\omega_0}\,\nu'(C'_\omega).
\]
Thus the refinement records only mass identities that were either supplied
atom-by-atom or supplied by the explicit finite-additivity hypotheses for the
two base cells.
\end{proof}
